凭借其性能，C++ 在软件行业的广大领域中保持着统治地位——编写良好的 C++ 代码几乎可以说比其它任何代码都跑得更快，因为 C++ 让程序员对编译器最终生成的代码拥有近乎完全的控制。

低级、高性能代码的经典特征之一，就是使用裸指针（\texttt{Foo*}）。然而，裸指针有许多陷阱，例如内存泄漏（memory leak）和悬垂（dangling）指针。C++11 标准库中的智能指针（smart pointer）类型可以帮助你以极小的代价（甚至毫无代价）避开这些陷阱。

本章我们将学习以下内容：

「智能指针」的定义，以及你如何自己动手写一个。

\texttt{std::\allowbreak{}unique\_\allowbreak{}ptr<\allowbreak{}T>\allowbreak{}} 在防止各类资源泄漏（不仅限于内存泄漏）方面的用处。

\texttt{std::\allowbreak{}shared\_\allowbreak{}ptr<\allowbreak{}T>\allowbreak{}} 是如何实现的，以及它对内存使用的影响。

奇异递归模板模式（CRTP）的含义与用途。

\section{智能指针的起源}

在 C 语言中，裸指针有许多用途：

作为调用方所拥有对象的一个廉价、不拷贝的视图（view）。

作为被调用方修改调用方所拥有对象的一种手段。

作为指针/长度对（pointer/length pair）的一半，用于数组。

作为一个可选参数（既可是一个有效指针，也可以是空指针）。

作为一种在堆上管理内存的手段。

在 C++ 中，我们有原生引用（\texttt{const Foo\&} 和 \texttt{Foo\&}）来处理前两项；此外，移动语义（move semantics）使得被调用方在大多数情况下按值获取并回传一个复杂对象的代价很低，从而完全避免了别名（aliasing）问题。在 C++17 中，我们可以使用 \texttt{std::\allowbreak{}string\_\allowbreak{}view} 来处理第一项和第三项中的部分情形。而我们在第 5 章「词汇类型」中刚刚看到，传一个 \texttt{optional<\allowbreak{}T>\allowbreak{}}——或者更花哨一点，用 \texttt{optional<\allowbreak{}reference\_\allowbreak{}wrapper<\allowbreak{}T>\allowbreak{}>}——就足以处理第四项。

本章要讨论的是第五项。

在 C 语言中，堆分配（heap allocation）带来一大堆问题，而所有这些问题（甚至更多！）在 2011 年之前的 C++ 中也同样存在。不过，到了 C++11，几乎所有这些问都消失了。我们来把它们列举一下：

内存泄漏（memory leak）：你可能在堆上分配了一块内存或一个对象，却不小心忘了写下释放它的代码。

内存泄漏（memory leak）：你可能写了那段代码，但由于提前返回（early return）或抛出了异常，代码从未运行，内存也就一直未被释放！

释放后使用（use-after-free）：你拷贝了一份指向堆上某个对象的指针，然后通过原始指针释放了该对象。持有那份拷贝指针的一方并不知道自己的指针已经失效。

通过指针算术破坏堆（heap corruption via pointer arithmetic）：你在地址 A 处的堆上分配了一个数组。手里拿着一个指向数组的裸指针，会诱使你做指针算术，最终你不小心释放了一个指向地址 A+k 的指针。当 k=0 时（墨菲定律保证在测试中必然如此）没有问题；当 k=1 时，你就破坏了堆并导致崩溃。

前两个问题因为堆分配在语义上允许失败而更加严重——\texttt{malloc} 可能返回空指针，\texttt{operator new} 可能抛出 \texttt{std::\allowbreak{}bad\_\allowbreak{}alloc}——这意味着，如果你在编写 C++11 之前会分配内存的代码，你很可能在写大量用于应对分配失败的清理代码。（在 C++ 中，无论你是否意识到，你都「在写」那些代码，因为由异常处理产生的控制流路径即便你没刻意去想，也依然存在。）这一切的结论是：在 C++ 中管理堆分配的内存是困难的。

除非你使用智能指针！

\section{智能指针永不遗忘}

「智能指针」类型（不要与「花式指针（fancy pointer）」类型混淆，后者我们将在第 8 章「分配器」中介绍）的思想是：它是一个类——通常是一个类模板——其句法行为就像一个指针，但它的特殊成员函数（构造、析构以及拷贝/移动）带有额外的记账（bookkeeping）逻辑，以确保某些不变式（invariant）。例如，我们可以确保以下几点：

指针的析构函数也会释放其所指向的对象——有助于解决内存泄漏问题。

也许该指针不可被拷贝——有助于解决释放后使用问题。

或者，该指针可以被拷贝，但它知道自己存在多少个副本，并且在指向它的最后一个指针被销毁之前不会释放所指对象。

或者，该指针可以被拷贝，你也可以释放所指对象，但如果你这么做了，所有指向它的其它指针会神奇地变成空指针。

或者，该指针没有内建的 \texttt{operator+}——有助于解决由指针算术导致的破坏问题。

或者，你可以对指针的值进行算术调整，但「指向哪个对象」的算术与「该释放哪个对象」的身份是分开管理的。

标准智能指针类型有 \texttt{std::\allowbreak{}unique\_\allowbreak{}ptr<\allowbreak{}T>\allowbreak{}}、\texttt{std::\allowbreak{}shared\_\allowbreak{}ptr<\allowbreak{}T>\allowbreak{}}，以及（严格说它不算指针类型，但我们把它也归在一起）\texttt{std::\allowbreak{}weak\_\allowbreak{}ptr<\allowbreak{}T>\allowbreak{}}。在本章中，我们将介绍这三种类型，外加一种你可能觉得有用的非标准智能指针——它未来说不定会成为某个 C++ 标准中的标准智能指针类型！

\section{用 \texttt{std::\allowbreak{}unique\_\allowbreak{}ptr<\allowbreak{}T>\allowbreak{}} 自动管理内存}

智能指针类型的根本属性很简单：它应该支持 \texttt{operator*}，并且应该重载特殊成员函数以维护其类不变式，无论这些不变式是什么。

\texttt{std::\allowbreak{}unique\_\allowbreak{}ptr<\allowbreak{}T>\allowbreak{}} 支持与 \texttt{T*} 相同的接口，但有一个类不变式：一旦你构造了一个指向某个堆分配对象的 \texttt{unique\_\allowbreak{}ptr}，当该 \texttt{unique\_\allowbreak{}ptr} 的析构函数被调用时，那个对象就会被释放。我们来写一些代码支持这个 \texttt{T*} 接口：

\begin{codebox}[cpp]
    template<typename T>
    class unique_ptr {
      T *m_ptr = nullptr;
    public:
      constexpr unique_ptr() noexcept = default;
      constexpr unique_ptr(T *p) noexcept : m_ptr(p) {}
      T *get() const noexcept { return m_ptr; }
      operator bool() const noexcept { return bool(get()); }
      T& operator*() const noexcept { return *get(); }
      T* operator->() const noexcept { return get(); }
\end{codebox}

如果我们到此为止——只有从 \texttt{T*} 构造指针对象的方法，以及把指针再取出来的方法——那我们就得到了本章末尾将讨论的 \texttt{observer\_\allowbreak{}ptr<\allowbreak{}T>\allowbreak{}}。但我们继续往下写。我们还会添加 \texttt{release} 和 \texttt{reset} 方法：

\begin{codebox}[cpp]
      void reset(T *p = nullptr) noexcept {
        T *old_p = std::exchange(m_ptr, p);
        delete old_p;
      }
      T *release() noexcept {
        return std::exchange(m_ptr, nullptr);
      }
\end{codebox}

\texttt{p.\allowbreak{}release(\allowbreak{})\allowbreak{}} 与 \texttt{p.\allowbreak{}get(\allowbreak{})\allowbreak{}} 很像，但除了返回原始裸指针的一份拷贝之外，它还会把 \texttt{p} 的内容置空（当然不会释放原始指针，因为大概我们的调用方想要取得它的所有权）。

\texttt{p.\allowbreak{}reset(\allowbreak{}q)\allowbreak{}} 会释放 \texttt{p} 当前的内容，然后把裸指针 \texttt{q} 放进去。

请注意，我们这两个成员函数都是借助标准算法 \texttt{std::\allowbreak{}exchange} 实现的，而该算法我们在第 3 章「迭代器对算法」中并未涉及。它有点像是 \texttt{std::\allowbreak{}swap} 的一个「返回值」版本：传入一个新值，取回旧值。

最后，有了这两个原语操作，我们就可以实现 \texttt{std::\allowbreak{}unique\_\allowbreak{}ptr<\allowbreak{}T>\allowbreak{}} 的特殊成员函数，从而维持我们的不变式——再次强调，这个不变式是：一旦一个裸指针被某个 \texttt{unique\_\allowbreak{}ptr} 对象获取，只要该 \texttt{unique\_\allowbreak{}ptr} 对象的值不变，它就一直有效；而一旦不再如此——当该 \texttt{unique\_\allowbreak{}ptr} 被调整去指向别处，或被销毁——那个裸指针就会被正确地释放。以下是这些特殊成员函数：

\begin{codebox}[cpp]
      unique_ptr(unique_ptr&& rhs) noexcept {
        this->reset(rhs.release());
      }
      unique_ptr& operator=(unique_ptr&& rhs) noexcept {
        reset(rhs.release());
        return *this;
      }
      ~unique_ptr() {
        reset();
      }
    };
\end{codebox}

在内存中，我们的 \texttt{std::\allowbreak{}unique\_\allowbreak{}ptr<\allowbreak{}T>\allowbreak{}} 看起来像这样：

\myGraphic{0.9}{images/fig_p157_1.png}{}

我们还需要一个小小的辅助函数，以便永远不必亲手去碰裸指针：

\begin{codebox}[cpp]
    template<typename T, typename... Args>
    unique_ptr<T> make_unique(Args&&... args)
    {
      return unique_ptr<T>(new T(std::forward<Args>(args)...));
    }
\end{codebox}

有了 \texttt{unique\_\allowbreak{}ptr} 这个工具，我们就可以替换掉下面这种老式代码：

\begin{codebox}[cpp]
    struct Widget {
      virtual ~Widget();
    };
    struct WidgetImpl : Widget {
      WidgetImpl(int size);
    };
    struct WidgetHolder {
      void take_ownership_of(Widget *) noexcept;
    };
    void use(WidgetHolder&);
    void test() {
      Widget *w = new WidgetImpl(30);
      WidgetHolder *wh;
      try {
        wh = new WidgetHolder();
      } catch (...) {
        delete w;
        throw;
      }
      wh->take_ownership_of(w);
      try {
        use(*wh);
      } catch (...) {
        delete wh;
        throw;
      }
      delete wh;
    }
\end{codebox}

它可以被替换成下面这样的现代 C++17 代码：

\begin{codebox}[cpp]
    void test() {
      auto w = std::make_unique<WidgetImpl>(30);
      auto wh = std::make_unique<WidgetHolder>();
      wh->take_ownership_of(w.release());
      use(*wh);
    }
\end{codebox}

请注意，\texttt{unique\_\allowbreak{}ptr<\allowbreak{}T>\allowbreak{}} 是 RAII（资源获取即初始化）的又一个应用——在这里，它名副其实。

虽然「有意思」的动作（底层裸指针的释放）仍然发生在（\texttt{unique\_\allowbreak{}ptr} 的）析构期间，但你要想获得 \texttt{unique\_\allowbreak{}ptr} 的全部收益，唯一的方式就是：确保每当你分配一个资源，你也同时初始化一个 \texttt{unique\_\allowbreak{}ptr} 来管理它。上一节展示的、并在 C++14 中加入标准库的 \texttt{std::\allowbreak{}make\_\allowbreak{}unique<\allowbreak{}T>\allowbreak{}(\allowbreak{})\allowbreak{}} 函数，就是现代 C++ 中安全内存管理的关键。

虽然不使用 \texttt{make\_\allowbreak{}unique} 也能用 \texttt{unique\_\allowbreak{}ptr}，但你不应该这么做：

\begin{codebox}[cpp]
    std::unique_ptr<Widget> bad(new WidgetImpl(30));
    bad.reset(new WidgetImpl(40));
    std::unique_ptr<Widget> good = std::make_unique<WidgetImpl>(30);
    good = std::make_unique<WidgetImpl>(40);
\end{codebox}

\subsection{为什么 C++ 没有 finally 关键字}

再考虑一下上一节「前现代」代码示例中的这段片段：

\begin{codebox}[cpp]
    try {
      use(*wh);
    } catch (...) {
      delete wh;
      throw;
    }
    delete wh;
\end{codebox}

在 Java、Python 等其它语言中，这些语义可以用 finally 关键字（finally keyword）表达得更紧凑：

\begin{codebox}[cpp]
    try {
      use(*wh);
    } finally {
      delete wh;
    }
\end{codebox}

C++ 没有 finally 关键字，而且毫无迹象表明它未来会进入这门语言。这纯粹是因为 C++ 与那些其它语言之间存在哲学上的差异：C++ 的哲学是，如果你关心的是要强制维护某个不变式——比如「无论我们以何种方式离开这个代码块，这个指针都应当在块末被释放」——那么你根本就不该去写显式的代码，因为那样你总有机会写错，从而引入 bug。

如果你有某个想要强制维护的不变式，那么强制维护它的正确位置在于类型系统之中，借助构造函数、析构函数以及其它特殊成员函数——也就是 RAII 的工具。这样，你就能确保对你这个新类型的任何可能使用都会维持其不变式——比如「只要底层指针不再被该类型的对象持有，它就应当被释放」——而当你去编写业务逻辑时，你将不必写任何显式代码；代码看起来简洁，却总能——可证明地——拥有正确的行为。

因此，如果你发现自己写的代码像前面的例子那样，或者你发现自己希望只要能写个 finally 就好了，请停下来想一想：「我是不是该为此用 \texttt{unique\_\allowbreak{}ptr}？」或者「我是不是该为此写一个 RAII 类类型？」

\subsection{自定义删除回调}

说到自定义的 RAII 类型，你或许会想知道：能否给 \texttt{std::\allowbreak{}unique\_\allowbreak{}ptr} 搭配一个自定义删除回调（deletion callback）？例如，与其把底层指针传给 \texttt{delete}，你可能想把它传给 \texttt{free(\allowbreak{})\allowbreak{}}。答案是可以的！

\texttt{std::\allowbreak{}unique\_\allowbreak{}ptr<\allowbreak{}T,\allowbreak{}D>\allowbreak{}} 有第二个模板类型参数：一个删除回调类型。参数 \texttt{D} 默认是 \texttt{std::\allowbreak{}default\_\allowbreak{}delete<\allowbreak{}T>\allowbreak{}}，它只是调用 \texttt{operator delete}，但你可以传入任何你想要的类型——通常是一个带有重载 \texttt{operator(\allowbreak{})\allowbreak{}} 的用户自定义类类型：

\begin{codebox}[cpp]
    struct fcloser {
      void operator()(FILE *fp) const {
        fclose(fp);
      }
      static auto open(const char *name, const char *mode) {
        return std::unique_ptr<FILE, fcloser>(fopen(name, mode));
      }
    };
    void use(FILE *);
    void test() {
      auto f = fcloser::open("test.txt", "r");
      use(f.get());
      // 即使 use() 抛出异常，f 也会被关闭
    }
\end{codebox}

顺带一提，请注意 \texttt{std::\allowbreak{}unique\_\allowbreak{}ptr} 的析构函数经过了精心编写，从而保证绝不会用一个空指针去调用你的回调。这在上面的例子中至关重要，因为 \texttt{fclose(\allowbreak{}NULL)\allowbreak{}} 是一个特例，含义是「关闭当前进程中的所有已打开文件句柄」——而这绝不是你想要做的事！

还要注意的是，\texttt{std::\allowbreak{}make\_\allowbreak{}unique<\allowbreak{}T>\allowbreak{}(\allowbreak{})\allowbreak{}} 永远只接受一个模板类型参数；并不存在 \texttt{std::\allowbreak{}make\_\allowbreak{}unique<\allowbreak{}T,\allowbreak{}D>\allowbreak{}(\allowbreak{})\allowbreak{}}。但「不要亲手去碰裸指针」这条规则依然是个好规则；正因如此，我们前面的例子把 \texttt{fopen} 和 \texttt{unique\_\allowbreak{}ptr} 的构造封装进了一个小巧的可复用辅助函数 \texttt{fcloser::\allowbreak{}open}，而不是把对 \texttt{fopen} 的调用内联进 \texttt{test} 的函数体。

你的自定义删除器（custom deleter）所需空间会被分配在 \texttt{std::\allowbreak{}unique\_\allowbreak{}ptr<\allowbreak{}T,\allowbreak{}D>\allowbreak{}} 对象自身内部，这意味着：如果 \texttt{D} 含有任何成员变量，那么 \texttt{sizeof(\allowbreak{}unique\_\allowbreak{}ptr<\allowbreak{}T,\allowbreak{}D>\allowbreak{})\allowbreak{}} 可能会大于 \texttt{sizeof(\allowbreak{}unique\_\allowbreak{}ptr<\allowbreak{}T>\allowbreak{})\allowbreak{}}：

\myGraphic{0.9}{images/fig_p161_1.png}{}

\subsection{用 \texttt{std::\allowbreak{}unique\_\allowbreak{}ptr<\allowbreak{}T[]>} 管理数组}

另一种 \texttt{delete p} 不是释放裸指针的恰当方式的情形是：如果 \texttt{p} 是一个指向数组首元素的指针；在这种情况下，应当改用 \texttt{delete [] p}。

所幸，从 C++14 起，\texttt{std::\allowbreak{}unique\_\allowbreak{}ptr<\allowbreak{}T[]>} 已经存在，并在这种情况下做正确的事（得益于 \texttt{std::\allowbreak{}default\_\allowbreak{}delete<\allowbreak{}T[]>} 同样存在并做正确的事，即调用 \texttt{operator delete[]}）。

确实存在针对数组类型的 \texttt{std::\allowbreak{}make\_\allowbreak{}unique} 重载，不过要小心——它给参数赋予了不同的含义！\texttt{std::\allowbreak{}make\_\allowbreak{}unique<\allowbreak{}T[]>(\allowbreak{}n)\allowbreak{}} 本质上调用 \texttt{new T[n](\allowbreak{})\allowbreak{}}，末尾的括号表示它会对所有元素进行值初始化（value-initialize）；也就是说，它会把基本类型清零。在极少数你不想如此的情形下，你就得自己调用 \texttt{new}，并把返回值包进一个 \texttt{std::\allowbreak{}unique\_\allowbreak{}ptr<\allowbreak{}T[]>}，最好像我们在上一节的例子（我们用的是 \texttt{fcloser::\allowbreak{}open}）那样使用一个辅助函数。

\section{用 \texttt{std::\allowbreak{}shared\_\allowbreak{}ptr<\allowbreak{}T>\allowbreak{}} 进行引用计数}

在彻底解决了内存泄漏问题之后，我们现在着手解决释放后使用（use-after-free）这类 bug。这里要解决的本质问题，是某个给定资源或内存块的所有权不清晰——或者更确切地说，是共享所有权（shared ownership）。这块内存可能在不同时刻有几个人在看着它，也许来自不同的数据结构，也许来自不同的线程，而我们想确保所有这些利益攸关方都参与到「何时释放它」的决策之中。底层那块内存的所有权应当是共享的。

为此，标准提供了 \texttt{std::\allowbreak{}shared\_\allowbreak{}ptr<\allowbreak{}T>\allowbreak{}}。它的接口看起来与 \texttt{std::\allowbreak{}unique\_\allowbreak{}ptr<\allowbreak{}T>\allowbreak{}} 非常相似；所有的差异都隐藏在幕后，也就是那些特殊成员函数的实现之中。

\texttt{std::\allowbreak{}shared\_\allowbreak{}ptr<\allowbreak{}T>\allowbreak{}} 提供了一种常被称为引用计数（reference counting）的内存管理方法。每一个由 \texttt{shared\_\allowbreak{}ptr} 管理的对象都记录着系统中有多少个指向它的引用——也就是此刻有多少个利益攸关方关心它——而一旦这个数字降为零，该对象就知道是时候清理自己了。当然，真正做清理的并不是「对象」本身；知道如何计数引用并清理东西的，是一个小巧的包装器，或称「控制块（control block）」，它会在你把某个对象的所有权转移给 \texttt{shared\_\allowbreak{}ptr} 时在堆上被创建出来。控制块由标准库在幕后处理，但如果我们去看它在内存中的布局，它大概长这样：

\myGraphic{0.9}{images/fig_p163_1.png}{}

正如 \texttt{unique\_\allowbreak{}ptr} 有 \texttt{make\_\allowbreak{}unique} 一样，标准库也为 \texttt{shared\_\allowbreak{}ptr} 提供了 \texttt{make\_\allowbreak{}shared}，让你永远不必亲手去碰裸指针。使用 \texttt{std::\allowbreak{}make\_\allowbreak{}shared<\allowbreak{}T>\allowbreak{}(\allowbreak{}args)\allowbreak{}} 来分配共享对象的另一个好处是：把所有权转入一个 \texttt{shared\_\allowbreak{}ptr} 需要为控制块分配额外的内存。当你调用 \texttt{make\_\allowbreak{}shared} 时，标准库被允许一次性分配一整块足以同时容纳控制块和你的 \texttt{T} 对象的内存。（这一点由前图中 \texttt{control\_\allowbreak{}block\_\allowbreak{}impl} 和 \texttt{Super} 这两个矩形的物理摆放位置所体现。）

拷贝一个 \texttt{shared\_\allowbreak{}ptr} 会使其关联控制块的「使用计数（use-count）」加一；销毁一个 \texttt{shared\_\allowbreak{}ptr} 会使其使用计数减一。对 \texttt{shared\_\allowbreak{}ptr} 赋值会使其旧值（若有）的使用计数减一，并使其新值的使用计数加一。下面是一些摆弄 \texttt{shared\_\allowbreak{}ptr} 的例子：

\begin{codebox}[cpp]
    std::shared_ptr<X> pa, pb, pc;
    pa = std::make_shared<X>();
    // 引用计数总是从 1 开始
    pb = pa;
    // 拷贝该指针；引用计数现在变为 2
    pc = std::move(pa);
    assert(pa == nullptr);
    // 移动该指针会使引用计数保持在 2
    pb = nullptr;
    // 将引用计数减回到 1
    assert(pc.use_count() == 1);
\end{codebox}

下面这张图说明了……的一个有趣且偶尔有用的方面。

\texttt{shared\_\allowbreak{}ptr}：两个 \texttt{shared\_\allowbreak{}ptr} 实例能够指向同一个控制块，却分别指向该控制块所管理的不同内存片段：

\myGraphic{0.9}{images/fig_p164_1.png}{}

上图中使用的这个构造函数（在 \texttt{get\_\allowbreak{}second(\allowbreak{})\allowbreak{}} 函数中也用到了它）通常被称为 \texttt{shared\_\allowbreak{}ptr} 的「别名构造函数」（aliasing constructor）。它接受一个任意类型的、非空的已有 \texttt{shared\_\allowbreak{}ptr} 对象，新构造出来的对象将与该对象共享同一个控制块。在下面的代码例子中，「destroying \texttt{Super}」这条消息要等到「accessing \texttt{Super::\allowbreak{}second}」这条消息打印之后才会输出：

\begin{codebox}[cpp]
    struct Super {
      int first, second;
      Super(int a, int b) : first(a), second(b) {}
      ~Super() { puts("destroying Super"); }
    };
    auto get_second() {
      auto p = std::make_shared<Super>(4, 2);
      return std::shared_ptr<int>(p, &p->second);
    }
    void test() {
      std::shared_ptr<int> q = get_second();
      puts("accessing Super::second");
      assert(*q == 2);
    }
\end{codebox}

如你所见，一旦所有权被移交给 \texttt{shared\_\allowbreak{}ptr} 系统，如何释放被管理资源的职责就完全落到了控制块的肩上。你的代码并不需要处处使用 \texttt{shared\_\allowbreak{}ptr<\allowbreak{}T>\allowbreak{}}，仅仅因为底层被管理的对象恰好是 \texttt{T} 类型。

\subsection{不要重复管理！}

虽然 \texttt{shared\_\allowbreak{}ptr<\allowbreak{}T>\allowbreak{}} 有潜力消除指针代码中那些讨厌的双重释放（double-free）缺陷，但令人遗憾地，缺乏经验的程序员太容易写出同样的缺陷——他们只是过度使用了那些接受裸指针参数的构造函数。下面是一个例子：

\begin{codebox}[cpp]
    std::shared_ptr<X> pa, pb, pc;
    pa = std::make_shared<X>();
      // 引用计数始终从 1 开始
    pb = pa;
      // 拷贝该指针；引用计数变为 2
    pc = std::shared_ptr<X>(pb.get()); // 错误！
      // 再次把同一个指针交给 shared_ptr，
      // 这会让 shared_ptr 对它进行两次管理！
    assert(pb.use_count() == 2);
    assert(pc.use_count() == 1);
    pc = nullptr;
      // pc 的引用计数降为 0，shared_ptr
      // 会对 X 对象调用 delete
    *pb; // 访问已释放的对象会导致未定义行为
\end{codebox}

请记住，你的目标应该是永远不要亲手去碰裸指针！这段代码出错的地方，正是它第一次调用 \texttt{pb.\allowbreak{}get(\allowbreak{})\allowbreak{}} 从 \texttt{shared\_\allowbreak{}ptr} 中取出裸指针的那一刻。

在这里调用别名构造函数 \texttt{pc = std::\allowbreak{}shared\_\allowbreak{}ptr<\allowbreak{}X>\allowbreak{}(\allowbreak{}pb,\allowbreak{} pb.\allowbreak{}get(\allowbreak{})\allowbreak{})\allowbreak{}} 本是正确的，但它的效果和简单的赋值 \texttt{pc = pb} 一样。所以我们可以得出另一条通用规则：如果你不得不在代码中显式写出 \texttt{shared\_\allowbreak{}ptr} 这个名字，那你多半是在做一些不寻常、甚至可能危险的事情。只要不在代码里写出 \texttt{shared\_\allowbreak{}ptr}，你已经可以分配并管理堆对象（使用 \texttt{std::\allowbreak{}make\_\allowbreak{}shared}），并且可以通过创建和销毁指针的拷贝（用 \texttt{auto} 按需声明变量）来操纵被管理对象的引用计数。这条规则唯一确定会失效的地方，是当你有时需要声明一个类型为 \texttt{shared\_\allowbreak{}ptr<\allowbreak{}T>\allowbreak{}} 的类数据成员时——这种时候一般不可能不写出类型名！

\subsection{用 \texttt{weak\_\allowbreak{}ptr} 持有可空句柄}

你可能已经在前面的图中注意到，控制块里有一个标记为「weak count」（弱计数）的、未曾解释的数据成员。现在是时候说明它是什么了。有时——这种情况很少见，但确实存在——我们会遇到这样的情形：我们用 \texttt{shared\_\allowbreak{}ptr} 来管理共享对象的所有权，却又想保留一个指向某对象的指针，而并不实际表达对该对象的所有权。当然，我们可以用裸指针、引用或 \texttt{observer\_\allowbreak{}ptr<\allowbreak{}T>\allowbreak{}} 来表达「非拥有引用」的概念，但那样的危险在于：被引用对象的实际所有者可能决定释放它，而当我们随后去解引用这个非拥有指针时，就会访问到一个已释放的对象，从而得到未定义行为（undefined behavior）。下面代码例子中的 \texttt{Dangerous\allowbreak{}Watcher} 就展示了这种危险的行为：

\begin{codebox}[cpp]
    struct DangerousWatcher {
      int *m_ptr = nullptr;
      void watch(const std::shared_ptr<int>& p) {
        m_ptr = p.get();
      }
      int current_value() const {
        // 此时，*m_ptr 可能已经被释放了！
        return *m_ptr;
      }
    };
\end{codebox}

我们也可以改用 \texttt{shared\_\allowbreak{}ptr} 来表达「引用」的概念，但那当然会给我们一个拥有式引用，使我们不再像 \texttt{Watcher}，而更像是 \texttt{Participant}：

\begin{codebox}[cpp]
    struct NotReallyAWatcher {
      std::shared_ptr<int> m_ptr;
      void watch(const std::shared_ptr<int>& p) {
        m_ptr = p;
      }
      int current_value() const {
        // 现在 *m_ptr 永远不会被释放；
        // 单是我们的存在就保住了 *m_ptr 的生命！
        return *m_ptr;
      }
    };
\end{codebox}

我们真正想要的是这样一种非拥有引用：它知晓 \texttt{shared\_\allowbreak{}ptr} 这套管理内存的机制，并且能查询控制块，判断被引用的对象是否依然存在。可是，当我们得知对象存在并准备去访问它时，它说不定已经被别的线程释放掉了！所以，我们需要的原语操作是「原子地获取一个指向被引用对象的拥有式引用（即一个 \texttt{shared\_\allowbreak{}ptr}）——如果对象存在的话，否则就表示失败」。也就是说，我们不想要一个非拥有引用；我们想要的是一张票据，可以在将来某个时刻兑换成一个拥有式引用。

标准库以 \texttt{std::\allowbreak{}weak\_\allowbreak{}ptr<\allowbreak{}T>\allowbreak{}} 这个名字提供了这种「\texttt{shared\_\allowbreak{}ptr} 的票据」。它之所以被称为「弱」（weak），正是相对于 \texttt{shared\_\allowbreak{}ptr} 的「强」拥有式引用而言的。下面是一个如何用 \texttt{weak\_\allowbreak{}ptr} 来解决我们 \texttt{Watcher} 问题的例子：

\begin{codebox}[cpp]
    struct CorrectWatcher {
      std::weak_ptr<int> m_ptr;
      void watch(const std::shared_ptr<int>& p) {
        m_ptr = std::weak_ptr<int>(p);
      }
      int current_value() const {
        // 现在我们可以安全地询问 *m_ptr 是否
        // 已被释放。
        if (auto p = m_ptr.lock()) {
            return *p;
        } else {
          throw "It has no value; it's been deallocated!";
        }
      }
    };
\end{codebox}

关于 \texttt{weak\_\allowbreak{}ptr}，你只需要知道两个操作：你可以从一个 \texttt{shared\_\allowbreak{}ptr<\allowbreak{}T>\allowbreak{}} 构造出 \texttt{weak\_\allowbreak{}ptr<\allowbreak{}T>\allowbreak{}}（通过调用构造函数，如上面的 \texttt{watch(\allowbreak{})\allowbreak{}} 函数所示）；你也可以通过调用 \texttt{wptr.\allowbreak{}lock(\allowbreak{})\allowbreak{}}，尝试从一个 \texttt{weak\_\allowbreak{}ptr<\allowbreak{}T>\allowbreak{}} 构造出 \texttt{shared\_\allowbreak{}ptr<\allowbreak{}T>\allowbreak{}}。如果 \texttt{weak\_\allowbreak{}ptr} 已经失效（expired），你会得到一个空的 \texttt{shared\_\allowbreak{}ptr}。

此外还有一个成员函数 \texttt{wptr.\allowbreak{}expired(\allowbreak{})\allowbreak{}}，它可以告诉你相关的 \texttt{weak\_\allowbreak{}ptr} 是否已经失效；但请注意，它基本上没什么用，因为即便它此刻返回 \texttt{false}，几微秒之后也可能返回 \texttt{true}。

下图承接前一幅图的故事：它从 \texttt{q} 创建了 \texttt{weak\_\allowbreak{}ptr}，随后把来源的那个 \texttt{shared\_\allowbreak{}ptr} 置空：

\myGraphic{0.9}{images/fig_p168_1.png}{}

拷贝一个 \texttt{weak\_\allowbreak{}ptr} 会使被引用对象控制块上的弱计数（weak-count）加一，而销毁一个 \texttt{weak\_\allowbreak{}ptr} 则使弱计数减一。当引用计数（use-count）降为零时，系统就知道可以安全地释放被控制的对象了；但控制块本身直到弱计数也降为零时才会被释放——到了那一刻，我们就知道再也没有任何 \texttt{weak\_\allowbreak{}ptr} 指向这个控制块了：

\myGraphic{0.9}{images/fig_p169_1.png}{}

你可能已经注意到，\texttt{shared\_\allowbreak{}ptr} 在它的 \texttt{Deleter} 上使用了与我们在第 5 章「词汇类型」中见到的 \texttt{std::\allowbreak{}any} 和 \texttt{std::\allowbreak{}function} 相同的技巧——它用的是类型擦除（type erasure）。和 \texttt{std::\allowbreak{}any}、\texttt{std::\allowbreak{}function} 一样，\texttt{std::\allowbreak{}shared\_\allowbreak{}ptr} 也提供了一个「碰碰运气」式的函数——\texttt{std::\allowbreak{}get\_\allowbreak{}deleter<\allowbreak{}Deleter>\allowbreak{}(\allowbreak{}p)\allowbreak{}}——用来取回原始的删除器（deleter）对象。这个小知识点在你的实际工作中完全没有用处；我提到它，只是为了提请你注意类型擦除在现代 C++ 中的重要性。即便是 \texttt{shared\_\allowbreak{}ptr} 这样表面目的与擦除类型毫无关系的组件，也在它功能的某个小角落里依赖于类型擦除。

\subsection{借助 \texttt{std::\allowbreak{}enable\_\allowbreak{}shared\_\allowbreak{}from\_\allowbreak{}this} 引用自身}

\texttt{shared\_\allowbreak{}ptr} 生态里还有最后一块内容值得我们讨论。此前已经提到过「双重管理」（double-managing）一个指针的危险——也就是创建多个控制块。因此，我们或许想要这样一种办法：给定一个指向堆分配对象的指针，能够确切地问出当前究竟是谁在管理它。

这个功能最常见的用例出现在面向对象编程中：某个方法 \texttt{A::\allowbreak{}foo(\allowbreak{})\allowbreak{}} 想要调用某个外部函数 \texttt{bar(\allowbreak{})\allowbreak{}}，而 \texttt{bar(\allowbreak{})\allowbreak{}} 需要一个指回 \texttt{A} 对象的指针。如果我们不考虑生命周期管理，这很容易做到——\texttt{A::\allowbreak{}foo(\allowbreak{})\allowbreak{}} 只需简单地调用 \texttt{bar(\allowbreak{}this)\allowbreak{}}。但假设我们的 \texttt{A} 正被 \texttt{shared\_\allowbreak{}ptr} 管理着，再假设 \texttt{bar(\allowbreak{})\allowbreak{}} 很可能会把 \texttt{this} 指针的一份拷贝藏在内部某处——也许是在注册一个稍后调用的回调，也许是在启动一个新线程，该线程会在 \texttt{A::\allowbreak{}foo(\allowbreak{})\allowbreak{}} 执行完毕并返回调用者之后并发运行。那么我们就需要某种办法，在 \texttt{bar(\allowbreak{})\allowbreak{}} 仍在运行期间保住 \texttt{A} 的生命。

显然，\texttt{bar(\allowbreak{})\allowbreak{}} 应该接受一个 \texttt{std::\allowbreak{}shared\_\allowbreak{}ptr<\allowbreak{}A>\allowbreak{}} 类型的参数，这样就能保住我们的 \texttt{A}。但是，在 \texttt{A::\allowbreak{}foo(\allowbreak{})\allowbreak{}} 内部，我们该从哪里得到那个 \texttt{shared\_\allowbreak{}ptr} 呢？我们可以给 \texttt{A} 一个 \texttt{std::\allowbreak{}shared\_\allowbreak{}ptr<\allowbreak{}A>\allowbreak{}} 类型的成员变量，可那样一来 \texttt{A} 就会保自己不死——它永远也不会消亡！这当然不是我们想要的结果！

一种初步的解决方案是：\texttt{A} 应该持有一个指向自身的、类型为 \texttt{std::\allowbreak{}weak\_\allowbreak{}ptr<\allowbreak{}A>\allowbreak{}} 的成员变量，并且在调用 \texttt{bar} 时，通过 \texttt{bar(\allowbreak{}this->m\_\allowbreak{}wptr.\allowbreak{}lock(\allowbreak{})\allowbreak{})\allowbreak{}} 来调用。不过这样做有相当多的语法开销，而且 \texttt{m\_\allowbreak{}wptr} 这个指针究竟该如何初始化也说不清楚。于是，C++ 把这个想法直接做进了标准库！

\begin{codebox}[cpp]
    template<class T>
    class enable_shared_from_this {
      weak_ptr<T> m_weak;
    public:
      enable_shared_from_this(const enable_shared_from_this&) {}
      enable_shared_from_this& operator=(const enable_shared_from_this&) {}
      shared_ptr<T> shared_from_this() const {
        return shared_ptr<T>(m_weak);
      }
    };
\end{codebox}

\texttt{std::\allowbreak{}enable\_\allowbreak{}shared\_\allowbreak{}from\_\allowbreak{}this<\allowbreak{}A>\allowbreak{}} 类持有我们那个类型为 \texttt{std::\allowbreak{}weak\_\allowbreak{}ptr<\allowbreak{}A>\allowbreak{}} 的成员变量，并以 \texttt{x.\allowbreak{}shared\_\allowbreak{}from\_\allowbreak{}this(\allowbreak{})\allowbreak{}} 这个名字暴露出「获取一个指向自身的 \texttt{shared\_\allowbreak{}ptr}」这一操作。在前面的代码中有几个值得注意的细节：第一，如果你试图在一个当前并未被 \texttt{shared\_\allowbreak{}ptr} 系统管理的对象上调用 \texttt{x.\allowbreak{}shared\_\allowbreak{}from\_\allowbreak{}this(\allowbreak{})\allowbreak{}}，你会得到一个 \texttt{std::\allowbreak{}bad\_\allowbreak{}weak\_\allowbreak{}ptr} 类型的异常。第二，请注意那个空的拷贝构造函数和拷贝赋值运算符。在这里，空的花括号和 \texttt{=default} 并不是一回事！如果我们用 \texttt{=default} 让拷贝操作取默认值，它们就会执行按成员拷贝。每当你拷贝一个被管理的对象，新对象就会收到原对象 \texttt{m\_\allowbreak{}weak} 的一份拷贝——而这根本不是我们想要的。一个 C++ 对象中 \texttt{enable\_\allowbreak{}shared\_\allowbreak{}from\_\allowbreak{}this} 部分的「身份」是和它所在的内存位置绑定在一起的，因此它不会（也不应该）遵循我们通常所追求的、关于拷贝与值语义（value semantics）的那些规则。

最后一个要回答的问题是：成员 \texttt{m\_\allowbreak{}weak}（请记住，它是一个私有成员；我们纯粹是为了讲解才使用 \texttt{m\_\allowbreak{}weak} 这个名字）究竟是如何被初始化的？答案是，\texttt{shared\_\allowbreak{}ptr} 的构造函数里包含一些代码，用来检测 \texttt{T} 是否公开继承自 \texttt{enable\_\allowbreak{}shared\_\allowbreak{}from\_\allowbreak{}this<\allowbreak{}T>\allowbreak{}}，如果是，就通过某个隐藏的后门来设置它的 \texttt{m\_\allowbreak{}weak} 成员。请注意，这种继承必须是公开且明确的，因为就 C++ 的规则而言，\texttt{shared\_\allowbreak{}ptr} 的构造函数不过是又一个用户定义的函数；它无法撬开你的类去寻找私有基类，也无法在 \texttt{enable\_\allowbreak{}shared\_\allowbreak{}from\_\allowbreak{}this} 的多个副本之间做消歧。

这些限制意味着：你应该永远只公开继承 \texttt{enable\_\allowbreak{}shared\_\allowbreak{}from\_\allowbreak{}this}；一旦某个类继承自 \texttt{enable\_\allowbreak{}shared\_\allowbreak{}from\_\allowbreak{}this}，你就应该永远只公开地继承它；而且，为了让事情简单，你可能应当只在继承体系的叶子节点上继承 \texttt{enable\_\allowbreak{}shared\_\allowbreak{}from\_\allowbreak{}this}。当然，如果你本来就不去构建深层继承体系，遵循这些规则就会相当容易！

让我们把关于 \texttt{enable\_\allowbreak{}shared\_\allowbreak{}from\_\allowbreak{}this} 的所有认识汇总到一个例子里：

\begin{codebox}[cpp]
    struct Widget : std::enable_shared_from_this<Widget> {
      template<class F>
      void call_on_me(const F& f) {
        f(this->shared_from_this());
      }
    };
    void test() {
      auto sa = std::make_shared<Widget>();
      assert(sa.use_count() == 1);
      sa->call_on_me([](auto sb) {
        assert(sb.use_count() == 2);
      });
      Widget w;
      try {
        w.call_on_me([](auto) {});
      } catch (const std::bad_weak_ptr&) {
        puts("Caught!");
      }
    }
\end{codebox}

\subsection{奇异递归模板模式}

你可能已经注意到了，不过在看过上面的代码例子之后应当更加明显：每当你继承 \texttt{enable\_\allowbreak{}shared\_\allowbreak{}from\_\allowbreak{}this} 时，你自己类的名字总会出现在它自身基类的模板参数列表中！这种「\texttt{X} 继承自 \texttt{A<\allowbreak{}X>\allowbreak{}}」的模式被称为\emph{奇异递归模板模式（Curiously Recurring Template Pattern，简称 CRTP）}。只要基类的某些方面依赖于它的派生类，这种模式就很常见。例如在我们的例子中，派生类的名字就被纳入了 \texttt{shared\_\allowbreak{}from\_\allowbreak{}this} 方法的返回类型中。

另一个适合使用 CRTP 的常见情形，是当派生类的某些行为被并入基类所提供的行为中时。例如，利用 CRTP，我们可以编写一个基类模板，为任何实现了 \texttt{operator+=} 和拷贝构造的派生类提供一个返回值的 \texttt{operator+}。请注意其中必需的 \texttt{static\_\allowbreak{}cast}：从 \texttt{addable<\allowbreak{}Derived>\allowbreak{}} 转换到 \texttt{Derived}，这样我们调用的就是 \texttt{Derived} 的拷贝构造函数，而不是基类 \texttt{addable<\allowbreak{}Derived>\allowbreak{}} 的拷贝构造函数：

\begin{codebox}[cpp]
    template<class Derived>
    class addable {
    public:
      auto operator+(const Derived& rhs) const {
        Derived lhs = static_cast<const Derived&>(*this);
        lhs += rhs;
        return lhs;
      }
    };
\end{codebox}

事实上，这几乎就是 Boost 的 Operators 库里 \texttt{boost::\allowbreak{}addable} 所提供的服务；不同之处在于，\texttt{boost::\allowbreak{}addable} 使用了所谓的「Barton-Nackman 技巧」，把它的 \texttt{operator+} 做成了一个友元自由函数，而不是成员函数：

\begin{codebox}[cpp]
    template<class Derived>
    class addable {
    public:
      friend auto operator+(Derived lhs, const Derived& rhs) {
        lhs += rhs;
        return lhs;
      }
    };
\end{codebox}

即便你从不在代码库里使用 \texttt{enable\_\allowbreak{}shared\_\allowbreak{}from\_\allowbreak{}this}，你也应该了解奇异递归模板模式（CRTP），并且能够在需要把某些派生类行为「注入」到基类方法里时，随时从你的工具箱中取出它。

\subsection{最后的警告}

\texttt{shared\_\allowbreak{}ptr}、\texttt{weak\_\allowbreak{}ptr} 和 \texttt{enable\_\allowbreak{}shared\_\allowbreak{}from\_\allowbreak{}this} 这套小生态是现代 C++ 中最酷的部分之一；它能让你的代码拥有垃圾回收语言般的安全性，同时又保留了速度，以及 C++ 一贯所具有的确定性析构。然而，要小心别滥用 \texttt{shared\_\allowbreak{}ptr}！你的大部分 C++ 代码根本不应该使用 \texttt{shared\_\allowbreak{}ptr}，因为你本不该去共享堆分配对象的所有权。你的第一选择应当永远是彻底避免堆分配（通过采用值语义）；第二选择应当是确保每个堆分配对象都有唯一的拥有者（通过使用 \texttt{std::\allowbreak{}unique\_\allowbreak{}ptr<\allowbreak{}T>\allowbreak{}}）；只有当这两者确实都做不到时，你才应该考虑使用共享所有权和 \texttt{std::\allowbreak{}shared\_\allowbreak{}ptr<\allowbreak{}T>\allowbreak{}}。

\section{用 \texttt{observer\_\allowbreak{}ptr<\allowbreak{}T>\allowbreak{}} 表示非特殊所有权}

到目前为止，我们已经见识了两种或三种不同的智能指针类型（取决于你是否把 \texttt{weak\_\allowbreak{}ptr} 算作一个独立的指针类型，还是更像是 \texttt{shared\_\allowbreak{}ptr} 的一张票据）。这些类型每一个都在源码层面携带了关于生命周期管理的一些有用信息。举个例子，仅凭下面这两个 C++ 函数的签名，我们能说出它们各自的语义吗？

\begin{codebox}[cpp]
    void remusnoc(std::unique_ptr<Widget> p);
    std::unique_ptr<Widget> recudorp();
\end{codebox}

我们看到，\texttt{remusnoc} 按值接受一个 \texttt{unique\_\allowbreak{}ptr}，这意味着被控制对象的所有权被转移给了 \texttt{remusnoc}。当我们调用这个函数时，我们必须对一个 \texttt{Widget} 拥有唯一所有权；而在调用之后，我们将再也无法访问那个 \texttt{Widget}。我们不知道 \texttt{remusnoc} 究竟是要销毁这个 \texttt{Widget}、把它保留下来，还是把它挂到别的什么对象或线程上；但可以明确的是，这已经不再是我们该操心的事了。\texttt{remusnoc} 函数是一个 widget 的消费者。

更微妙的是，我们还可以说：当我们调用 \texttt{remusnoc} 时，我们必然拥有一个用 \texttt{new} 分配、且可以安全 \texttt{delete} 的 \texttt{Widget} 的唯一所有权！

反过来也一样：当我调用 \texttt{recudorp} 时，我知道收到的 \texttt{Widget} 一定由我独占拥有。它不是对别人持有的 \texttt{Widget} 的引用，也不是指向某块静态数据的指针。它明确是一个由我、且仅由我拥有的、分配在堆上的 \texttt{Widget}。即便我拿到返回值后的第一件事，就是对其调用 \texttt{.release(\allowbreak{})\allowbreak{}}，并把那个裸指针塞进某个“前现代”结构体里，我也能确信这样做是安全的，因为我毫无疑问是这个返回值的唯一拥有者。

关于这个 C++ 函数的语义，我们能说些什么呢？

\begin{codebox}[cpp]
    void suougibma(Widget *p);
\end{codebox}

这很含糊。这个函数也许会接管传入指针的所有权，也许不会。我们只能（但愿如此）从 \texttt{suougibma} 的文档里，或者从代码库的某些风格约定（比如“裸指针绝不表示所有权”——这是个合理的约定）里去判断，但仅凭函数签名无法断定。表达这种区别的另一种说法是：\texttt{unique\_\allowbreak{}ptr<\allowbreak{}T>\allowbreak{}} 是用来表达所有权转移的词汇类型，而 \texttt{T*} 根本算不上任何东西的词汇类型；它相当于 C++ 里的无意义词或罗夏墨迹，因为不同的人未必会对它的含义达成一致。

所以，如果你发现自己的代码库里到处都在传递非拥有型指针，你大概会想要一个词汇类型来表示“非拥有型指针”这个概念。（你的第一步应当是尽可能用引用而非指针来传参，不过这里我们假设这条路你已经走不通了。）这样的词汇类型确实存在，尽管它还（暂时）不在 C++ 标准库里：它出自 Walter Brown 之手，被称为“世界上最笨的智能指针”，本质上只是把一个非拥有型裸指针用类的形式包了一层：

\begin{codebox}[cpp]
    template<typename T>
    class observer_ptr {
      T *m_ptr = nullptr;
      public:
      constexpr observer_ptr() noexcept = default;
      constexpr observer_ptr(T *p) noexcept : m_ptr(p) {}
      T *get() const noexcept { return m_ptr; }
      operator bool() const noexcept { return bool(get()); }
      T& operator*() const noexcept { return *get(); }
      T* operator->() const noexcept { return get(); }
    };
    void revresbo(observer_ptr<Widget> p);
\end{codebox}

有了 \texttt{observer\_\allowbreak{}ptr} 这个工具，\texttt{revresbo} 只是观察其参数就变得一目了然；它显然不会接管该参数的所有权。事实上，我们可以假设它甚至不会保存传入指针的副本，因为那个指针的有效性取决于被控对象的生存期，而 \texttt{revresbo} 明确声明自己对该对象的生存期没有任何利害关系。如果它想对被控对象的生存期拥有利害关系，它就会向调用者显式索取——要求传入 \texttt{unique\_\allowbreak{}ptr} 或 \texttt{shared\_\allowbreak{}ptr}。通过要求 \texttt{observer\_\allowbreak{}ptr}，\texttt{revresbo} 等于在整场所有权之争中置身事外了。

如前所述，\texttt{observer\_\allowbreak{}ptr} 并不属于 C++17 标准。阻碍它进入标准的主要反对意见之一，就是它那个糟糕的名字（因为它与“观察者模式”毫无关系）。另外也有许多行家认为，\texttt{T*} 才应当是表示“非拥有型指针”的词汇类型，所有用 \texttt{T*} 来做所有权转移的老代码，都应当被重写，或至少用 \texttt{owner<\allowbreak{}T*>} 之类的构造重新标注。这正是《C++ 核心指南》（C++ Core Guidelines）的编辑者们——包括 C++ 之父 Bjarne Stroustrup——目前推荐的做法。不过有一点是确凿无疑的：绝不要用裸指针来做所有权转移！

\section{小结}

本章中，我们学到了关于智能指针的一些知识。

\texttt{std::\allowbreak{}unique\_\allowbreak{}ptr<\allowbreak{}T>\allowbreak{}} 是表达所有权以及所有权转移的词汇类型；应优先使用它，而非裸 \texttt{T*}。在那些明确不发生所有权转移的情形下，或裸 \texttt{T*} 可能让读者产生歧义的情形下，可以考虑使用 \texttt{observer\_\allowbreak{}ptr}。

\texttt{std::\allowbreak{}shared\_\allowbreak{}ptr<\allowbreak{}T>\allowbreak{}} 是处理共享所有权的好工具（而且是标准工具），所谓共享所有权，是指多个不同的实体都是同一个被控对象生存期的利益相关方。\texttt{std::\allowbreak{}weak\_\allowbreak{}ptr<\allowbreak{}T>\allowbreak{}} 是 \texttt{shared\_\allowbreak{}ptr} 的“票券”；它提供 \texttt{.lock(\allowbreak{})\allowbreak{}} 而非 \texttt{operator*}。如果你的类需要获取指向自身的 \texttt{shared\_\allowbreak{}ptr}，就让它继承 \texttt{std::\allowbreak{}enable\_\allowbreak{}shared\_\allowbreak{}from\_\allowbreak{}this<\allowbreak{}T>\allowbreak{}}。记住要公有继承，并且一般来说，只在继承图的叶节点上这样继承。此外，在并非绝对需要共享所有权的情形下，不要滥用这些特性！

绝不要亲手去碰裸指针：用 \texttt{make\_\allowbreak{}unique} 和 \texttt{make\_\allowbreak{}shared} 来创建分配在堆上的对象，并一举把它们管理起来。另外，每当你需要把派生类的行为“注入”回基类提供的函数中时，别忘了奇异递归模板模式（CRTP）。

在下一章，我们将讨论另一种“共享”：那种出现在多线程编程中的共享。
